\documentclass[11pt]{article}
% Frontiers-portable stand-in (article class). At submission, port to the
% Frontiers .cls; the sectioning, abstract, and keywords map directly.
\usepackage[margin=1in]{geometry}
\usepackage{amsmath,amssymb}
\usepackage{graphicx}
\graphicspath{{../figures/}}
\usepackage{booktabs}
\usepackage[round]{natbib}
\usepackage{siunitx}
\usepackage{xcolor}
\usepackage[hidelinks]{hyperref}
\newcommand{\todo}[1]{\textcolor{red}{[#1]}}
\newcommand{\bp}{\ensuremath{\,\mathrm{bp}}}
\title{\textbf{A Physics-Informed Neural Network for Option Pricing under Rough
Volatility: Meshfree Solution of the Markovian-Lifted PDE and Fast Parametric
Calibration}}
\author{
Dr. Aviral Sharma\\
Assistant Professor, AIT, Chandigarh University
\and
Kamal Preet Singh\\
Assistant Professor, SGT University
\and
Dr. Vipul Sharma\\
Independent Researcher
}

\begin{document}
\maketitle

\begin{abstract}
Rough-volatility models capture how equity and cryptocurrency option markets
actually behave, but they carry a computational curse: the driving noise is
fractional, the variance process is non-Markovian, and there is no
finite-dimensional pricing partial differential equation (PDE) to solve. The
Markovian lift buys a PDE back by approximating the fractional kernel with $n$
exponential factors---at the price of turning a two-dimensional problem into an
$(n{+}1)$-dimensional one. We take this on with a physics-informed neural network
(PINN) that learns the price by minimizing the strong-form PDE residual, and we
ask a sharp question: as the lifted dimension grows, does the network's own
accuracy collapse the way a finite-difference grid would? It does not. Measuring
the network's \emph{solve-error} against an independent lifted Monte Carlo at the
same $n$, and separating it cleanly---in price space---from the
model-approximation \emph{lift-error}, we find the solve-error grows only
sub-linearly, staying under about eleven price basis points out to $n{=}32$,
whereas a naive tensor grid would need $50^{n+1}$ points and is already infeasible
by $n{=}8$. The crossover between the two error sources gives a practical rule for
choosing $n$, and we show it shifts with maturity. Calibrated to real Bitcoin
options, the model returns a genuinely rough Hurst exponent near $0.09$ with
strong negative skew. Finally we build a parametric PINN that takes the model
parameters as inputs, holding the spot variance and maturity fixed; on its own it
is too coarse for production calibration, but---and this is the payoff---using it
to warm-start a handful of true-model polish steps recovers full from-scratch
calibration quality at roughly a ninth of the cost. The takeaway is that where
grid methods succumb to dimension, a mesh-free PINN degrades gracefully, and its
greatest practical value is as a differentiable warm-start for calibration.
\end{abstract}

\noindent\textbf{Keywords:} physics-informed neural networks; rough volatility;
option pricing; Markovian lift; curse of dimensionality; model calibration;
scientific machine learning.

% =====================================================================
\section{Introduction}

Every derivative-pricing problem is, underneath, a PDE. For the workhorse models
of quantitative finance that PDE is low-dimensional and a finite-difference grid
dispatches it in milliseconds. Rough volatility breaks this comfortable picture.
Driven by the empirical roughness of realized variance and by the steep
short-maturity skew that classical models such as \citet{heston1993closed} cannot
reproduce \citep{gatheral2018volatility,bennedsen2022decoupling,bayer2016pricing},
and with microstructural foundations in high-frequency trading
\citep{eleuch2018microstructural}, its variance is driven by a fractional noise
with Hurst exponent $H<1/2$. That single fact---an idea whose lineage runs back to
long-memory volatility \citep{comte1998long}---makes the variance neither
Markovian nor a semimartingale, and leaves \emph{no} finite-dimensional PDE to
discretize at all \citep{gatheral2018volatility,eleuch2019characteristic}. The
standard escape is the Markovian lift: approximate the fractional kernel by a sum
of $n$ exponentials and recover a genuine $(n{+}1)$-dimensional PDE
\citep{abijaber2019multifactor,abijaber2019lifting}. But the escape has a cost. It
trades an intractable problem for a high-dimensional one, and the accuracy knob
$n$ sits directly on top of the PDE dimension. That coupling is exactly what makes
the lifted rough-volatility PDE such a clean stress test for mesh-free neural
solvers and for the curse of dimensionality---and it is the test we run here.

Physics-informed neural networks are built for this setting. Rather than laying
down a grid, a PINN represents the price as a neural network and trains it to
satisfy the PDE at randomly sampled collocation points \citep{raissi2019physics}.
Nothing in the method scales with a grid, so---like the deep Galerkin and deep
BSDE solvers before it \citep{sirignano2018dgm,han2018solving}---it sidesteps the
exponential blow-up that kills finite differences in high dimensions. For pricing
it brings two further advantages we lean on: the physics is known exactly, so no
labelled prices are needed; and the trained network is differentiable in its
inputs, which we later exploit to calibrate by gradient descent. We run every
experiment on Bitcoin (BTC) options, a liquid market with a pronounced negative
skew that, as we show, calibrates to a firmly rough Hurst exponent.

We state our contributions exactly as strongly as the evidence allows, and no
stronger.
\begin{enumerate}
\item \textbf{A dimension-scaling result (our headline).} The network's
\emph{solve-error}---how well it actually solves the PDE it is handed---grows only
\emph{sub-linearly} as the lifted dimension climbs, staying below roughly eleven
price basis points out to $n=32$, while a naive full-tensor grid explodes as
$50^{n+1}$ and is infeasible by $n=8$. We are careful about the comparison:
$50^{n+1}$ is the \emph{naive} grid, and structured schemes do better but still
degrade with $n$, so our claim is that finite differences here are
\emph{impractical}, not impossible.
\item \textbf{An error decomposition and a rule for $n$.} We split the error into
the part the network controls (solve-error, versus a lifted Monte Carlo at the
same $n$) and the part only more factors can fix (lift-error, versus the true
model), report them in price space where they add exactly, and turn their
crossover into a concrete rule for choosing $n$ that shifts with maturity.
\item \textbf{A real calibration.} Fit to the live Deribit BTC surface, rough
Heston returns a rough Hurst exponent with strong negative correlation.
\item \textbf{Fast recalibration, honestly.} A parametric PINN plus a warm-start
hybrid recovers from-scratch calibration quality at about nine times fewer
true-model evaluations---and we are candid that the network on its own does not
get there.
\end{enumerate}

We are also clear about what is \emph{not} new. We are not the first to solve the
lifted (rough) Heston PDE with a mesh-free network: \citet{papapantoleon2025timestepping}
got there first, by a different route. What we add is the combination---a residual
PINN, a quantitative scaling-and-decomposition analysis, a calibration to real
crypto data, and the recalibration hybrid---which we develop next.

% =====================================================================
\section{Related work}

\paragraph{Mesh-free deep pricing on the lifted PDE (the closest work).}
\citet{papapantoleon2025timestepping} price European options in (rough) diffusion
models by recasting the pricing PDE as an energy-minimization problem and
approximating it in a time-stepping ``deep gradient flow,'' demonstrated on the
lifted Heston model and careful about the large-moneyness asymptotics and known
price bounds. This is the work nearest ours in both target (the lifted rough-vol
PDE) and spirit (mesh-free deep pricing), and we cite it precisely so as not to
overreach. It differs in method---theirs is a time-stepping variational scheme,
ours a PINN that minimizes the strong-form residual with a hard-wired initial
condition---and in scope: our dimension-scaling law, solve/lift decomposition,
market calibration, and recalibration hybrid are not its concerns.

\paragraph{Other neural solvers for rough volatility.}
\citet{jacquier2023random} attack the same model class from a very different
angle, casting the path-dependent pricing equations as backward stochastic
differential equations and solving them with a reservoir-type \emph{random}
network, which collapses training to a least-squares regression with convergence
guarantees. We share the target but not the mechanism, and neither dimension
scaling nor calibration speed is studied there.

\paragraph{Neural calibration surrogates (the foil).}
The dominant recipe for fast (rough-)volatility calibration is the supervised
surrogate of \citet{bayer2019deep} and \citet{horvath2021deep}: train a network
offline on a large grid of precomputed prices or implied volatilities, then use it
as a fast pricing map inside a calibration loop. It is purely data-driven and
carries no PDE consistency---the network never sees the governing equation. Ours
is \emph{physics-informed}: it learns from the residual, not from precomputed
labels. That distinction is what makes our hybrid work---the physics-trained
network supplies a warm start that a few true-model evaluations then sharpen
(Section~\ref{sec:hybrid}).

\paragraph{Foundations.} We build on PINNs \citep{raissi2019physics} and, more
broadly, mesh-free deep solvers for high-dimensional PDEs
\citep{sirignano2018dgm,han2018solving}; on rough volatility
\citep{gatheral2018volatility,eleuch2019characteristic} and its multifactor
Markovian lift \citep{abijaber2019multifactor,abijaber2019lifting}; and on
classical transform pricing \citep{heston1993closed,carr1999option} for our
independent ground truth. Two training ingredients are worth flagging up front,
because PINNs are notorious for optimization pathologies
\citep{krishnapriyan2021characterizing,wang2022when}: self-adaptive loss weighting
by homoscedastic uncertainty \citep{kendall2018multitask,mcclenny2023self}, and
causal (curriculum) training in time \citep{wang2024respecting}.

% =====================================================================
\section{Methods}

\subsection{Rough Heston and the Markovian lift}
Under rough Heston the variance is driven by the fractional kernel
$K(t)=t^{\alpha-1}/\Gamma(\alpha)$ with $\alpha=H+\tfrac12$ and $H<\tfrac12$,
which is non-Markovian \citep{gatheral2018volatility,eleuch2019characteristic}.
The lift \citep{abijaber2019multifactor,abijaber2019lifting} replaces it with
$K(t)\approx\sum_{i=1}^{n}c_i e^{-x_i t}$, introducing $n$ Ornstein--Uhlenbeck
factors $U^i$ with $V=V_0+\sum_i c_i U^i$. The discounted price $P(t,x,u)$, with
$x=\log S$, then solves the $(n{+}1)$-dimensional backward Kolmogorov PDE
\begin{align}
0 = {}& P_t + \Big(r-\tfrac12 V\Big)P_x + \tfrac12 V\,P_{xx}
   + \sum_i\big[-x_i u_i + \lambda(\theta-V)\big]P_{u_i} \nonumber\\
   & + \tfrac12\nu^2 V\sum_{i,j}P_{u_i u_j}
   + \rho\nu V\sum_i P_{x u_i} - rP,
\label{eq:pde}
\end{align}
with terminal condition $P(T,x,u)=\mathrm{payoff}(e^{x})$. All factors share one
Brownian motion, which is why the full double sum and the cross term appear. The
knob $n$ therefore \emph{is} the PDE dimension. We build $(c_i,x_i)$ by the
geometric-quadrature rule of \citet{abijaber2019multifactor}; its relative $L^2$
kernel error falls monotonically from $25.6\%$ at $n=5$ to $0.39\%$ at $n=100$, so
raising $n$ genuinely buys model fidelity.

\subsection{Two independent ground truths}
We hold the network to two references that we trust independently.
The first is a semi-analytic \emph{Fourier} pricer: it solves the fractional
Riccati equation for the rough Heston characteristic function
\citep{eleuch2019characteristic} and inverts it by Fourier integration
\citep{carr1999option}. It never uses the lift, so it is the ground truth of the
\emph{true model} ($n\to\infty$); in the flat-volatility limit it reproduces
Black--Scholes to $3.7\times10^{-7}$. The second is a \emph{lifted Monte Carlo}
simulator that integrates the $n$-factor lifted SDE with an exponential-integrator
Euler scheme handling the stiff $-x_i U^i$ term exactly---without which the scheme
diverges, since the mean reversions span some ten orders of magnitude for small
$H$. It matches Black--Scholes within Monte Carlo error. Because the two
approximate different objects (an ODE versus the full dynamics), their agreement
is a real cross-check. The subtle point we exploit throughout: \emph{at finite $n$
the lifted MC is the ground truth of the PDE the network solves}, while Fourier is
the ground truth of the model---and conflating the two is a trap we return to in
Section~\ref{sec:decomp}.

\subsection{The physics-informed network}
We solve \eqref{eq:pde} with a PINN \citep{raissi2019physics}. The network is a
$\tanh$ multilayer perceptron of $(\tau,x,u)$ with $\tau=T-t$, so the terminal
condition becomes an initial condition at $\tau=0$. Rather than penalize that
condition softly, we bake it in through the ansatz
\begin{equation}
P(\tau,x,u)=P_{\mathrm{base}}(\tau,x)+\tau\,N_\theta(\tau,x,u),
\label{eq:ansatz}
\end{equation}
whose $\tau$ prefactor forces the learned correction $N_\theta$ to vanish at
$\tau=0$ automatically. We anchor the baseline $P_{\mathrm{base}}$ not at the spot
variance $V_0$ but at the model's expected \emph{integrated} variance
$w(\tau)=\int_0^\tau \mathbb{E}[V_s]\,ds$, which already captures the mean
reversion of $V$ from $V_0$ toward $\theta$; anchoring at $V_0$ leaves a
systematic level bias for the network to clean up, and we saw exactly that bias
before we fixed the anchor.

We add asymptotic boundary losses---$P\to0$ as $S\to0$ and
$P\to S-Ke^{-r\tau}$ as $S\to\infty$---resampled each iteration. Because PINNs so
readily converge to the wrong solution when loss terms sit on very different
scales or when temporal causality is ignored
\citep{krishnapriyan2021characterizing,wang2022when}, we add two stabilizers that
proved decisive. First, we balance the residual and the two boundary terms with
self-adaptive per-term weights, using the learnable homoscedastic-uncertainty
scheme of \citet{kendall2018multitask} in the self-adaptive spirit of
\citet{mcclenny2023self}, so we never hand-tune a weight. Second, a
$\tau$-curriculum trains short maturities first and grows the horizon, respecting
causality in time \citep{wang2024respecting}. One more ingredient mattered more
than we expected: we do not sample the factors $u$ in a fixed symmetric box, but
draw each from the mean and spread it actually reaches under the lifted dynamics,
measured by a short pre-flight Monte Carlo. Getting this wrong flattens the
smile; getting it right restores the skew.

\subsection{How we score the network, and one trap we avoid}
\label{sec:decomp}
We train the network at fixed strike $K=1$ and price as a function of spot, then
read off the fixed-spot smile $C(1,K)$ through the call's degree-one homogeneity,
\begin{equation}
C(1,K)=K\,C\!\left(1/K,\,1\right),
\label{eq:homog}
\end{equation}
which holds because the variance dynamics do not depend on the spot level.
Equation~\eqref{eq:homog} is exact, not an approximation, and we apply it
consistently in every comparison in this paper. We stress this because a protocol
ablation---dropping the rescale and comparing the mismatched options
directly---inflates the apparent error to about $2700$ vol\bp. That number is a
property of the broken comparison, not a residual error of our method; our
accuracy is always the decomposed quantity we now define. Rather than a single
``PINN-versus-Fourier'' figure, we split the error per strike, in price space
where the signed contributions add exactly,
\begin{equation}
\underbrace{(\mathrm{PINN}-\mathrm{Fourier})}_{\text{total}}
=\underbrace{(\mathrm{PINN}-\mathrm{MC})}_{\text{solve-error}}
+\underbrace{(\mathrm{MC}-\mathrm{Fourier})}_{\text{lift-error}} .
\label{eq:split}
\end{equation}
The solve-error is ours to improve through training, architecture, and sampling;
the lift-error is a model approximation the network cannot touch and only $n$ can
shrink. We report price RMSE in basis points of spot, where \eqref{eq:split} is
additive, and quote implied-vol equivalents where they aid intuition; the two
components carry opposite signs, so they do not add in RMS.

\subsection{The parametric network}
For instant recalibration we train a single network that also takes the five model
parameters $(H,\nu,\rho,\lambda,\theta)$ as inputs, over a box around the
calibrated BTC values ($H\in[0.04,0.20]$, $\nu\in[0.10,0.60]$,
$\rho\in[-0.95,-0.20]$, $\lambda\in[0.5,4.0]$, $\theta\in[0.08,0.35]$), holding the
spot variance $V_0$ and the maturity fixed. Since the lift weights depend on $H$,
we precompute them on an $H$-grid and gather per collocation point while the
parameters enter the network as exact continuous inputs; we fix the factor count
by the crossover rule of Section~\ref{sec:nconv} ($n=10$). Crucially, the network
is differentiable in the model parameters, which is what lets us calibrate by
gradient descent.

\subsection{Calibration procedures and data}
We calibrate three ways and compare them head to head: (a)~\emph{Fourier-in-the-loop},
differential evolution with the Fourier pricer inside the objective;
(b)~\emph{PINN gradient}, Adam on the parameters using the parametric network as a
differentiable pricer, with zero Fourier calls; and (c)~\emph{hybrid}, the PINN's
single-shot solution as an initial guess, sharpened by a few local Fourier polish
steps (Nelder--Mead). The objective throughout is the vega-weighted
implied-volatility RMSE against the market smile. The market data are end-of-day
BTC option quotes from Deribit, a cryptocurrency derivatives exchange and data
provider, retrieved through its public interface; we set $r=0$ and infer the
forward from the quotes, as is standard for crypto, and recompute implied
volatilities from mid prices with a robust Black--Scholes inverter.

% =====================================================================
\section{Results}

\subsection{The network scales where the grid cannot}
\label{sec:nconv}
Our headline experiment asks what happens to the network's own accuracy as we push
the lifted dimension. We train one architecture (width $64$, depth $4$) at
$n=4,8,16,32$ on the canonical BTC calibration at $T=0.15$ and measure the
solve-error against the lifted MC over three seeds (Table~\ref{tab:nconv},
Figure~\ref{fig:nconv}). The answer is encouraging: the solve-error creeps up only
mildly, from $5.4\pm1.0$ to $10.7\pm2.0$ price\bp{}---about a doubling for an
eight-fold jump in dimension. We resist the temptation to call it flat; an early
single-seed run made it look flat, and three seeds corrected us. With only three
seeds the error bars at adjacent $n$ overlap, so we read the trend across the full
range rather than claiming significance between neighbours, and the robust
statement is that two-fold growth from $n=4$ to $n=32$. Wall-clock to the plateau
stays bounded throughout (about $6$--$10$ minutes on a CPU, with per-iteration
cost essentially flat at $55$--$62\,$ms) because the factor Hessian is vectorized.

Set this against the alternative. A naive full-tensor finite-difference grid on
the same $(n{+}1)$-dimensional PDE, at a modest $50$ nodes per dimension, needs
$50^{n+1}$ points: $3.1\times10^{8}$ at $n=4$, $2.0\times10^{15}$ at $n=8$, and
$1.2\times10^{56}$ at $n=32$. Even $n=8$---about $16$ petabytes in double
precision---is out of reach, while the network trains every $n$ on a laptop. We
are deliberately fair about this contrast: $50^{n+1}$ is the \emph{naive} grid, and
structured schemes (sparse grids, ADI, low-rank tensor-train) reach further, but
their cost and accuracy still degrade with $n$---sparse-grid node counts grow
super-linearly, and tensor-train needs a low-rank structure that the coupled
$\rho$--$\nu$ cross terms erode---and none is standard or validated for this PDE.
So the honest verdict is that a competitive finite-difference solve at
$n=16$--$32$ is \emph{impractical}, not provably impossible. Even in that
conservative form, gentle solve-error growth with bounded cost, against a grid
that cannot be stored, is precisely the curse-of-dimensionality advantage that
motivates going mesh-free.

\begin{table}[t]\centering
\caption{Dimension scaling (canonical BTC calibration, $T=0.15$; solve-error is
3-seed mean\,$\pm$\,SD; price RMSE in basis points of spot).}
\label{tab:nconv}
\begin{tabular}{ccccc}
\toprule
$n$ & solve-error (px\bp) & lift-error (px\bp) & train time (s) & FD grid $50^{(n+1)}$\\
\midrule
4  & $5.4\pm1.0$  & $16.4$ & $397$ & $3.1\times10^{8}$\\
8  & $6.1\pm1.4$  & $7.9$  & $386$ & $2.0\times10^{15}$\\
16 & $8.8\pm1.6$  & $3.5$  & $591$ & $7.6\times10^{28}$\\
32 & $10.7\pm2.0$ & $1.7$  & $457$ & $1.2\times10^{56}$\\
\bottomrule
\end{tabular}
\end{table}

\begin{figure}[t]\centering
\includegraphics[width=0.78\linewidth]{fig1_n_convergence.png}
\caption{Dimension scaling. The PINN solve-error (blue, 3-seed mean\,$\pm$\,SD)
grows sub-linearly and stays under $\sim\!11$ price\bp{} to $n=32$, while the
lift-error (red) falls monotonically and training time (green) stays bounded; the
naive finite-difference grid (inset) is infeasible by $n=8$.}
\label{fig:nconv}
\end{figure}

\subsection{Splitting the error, and a rule for how many factors to use}
The decomposition did more than tidy our bookkeeping---it overturned how we read
our own numbers. Taken at face value, a single ``PINN vs Fourier'' implied-vol
RMSE of $114$--$127$ vol\bp{} at $n=4$ looks like network error. Split by
\eqref{eq:split}, almost all of it is the lift. In the full configuration the
$n=4$ error is $37$ vol\bp{} of solve-error on top of $133$ of lift (they partly
cancel, to a total of $121$), and with correlation switched off it is a mere $8$
of solve-error under $133$ of lift. In other words, the network already solves its
PDE to well under $50$ vol\bp{}; what had looked like a modelling failure of the
network was the finite-$n$ model catching up to reality.

Because the lift-error shrinks monotonically with $n$ ($133,63,27,12$ vol\bp{} at
$n=4,8,16,32$; equivalently $16.4\to1.7$ price\bp) while the solve-error climbs
gently, the two curves cross. At $T\approx0.15$ they meet near $n=8$--$10$
(Figure~\ref{fig:nconv}), and that crossing is a genuinely useful design rule:
below it, add factors, because model approximation dominates; above it, stop,
because you are only paying for network error. The crossover is not fixed---it
tracks maturity. Over a $T$-sweep the solve-error stays small and bounded
everywhere (roughly $1$--$9.5$ price\bp{} across $T\in\{0.05,0.15,0.5\}$ and
$H\in\{0.05,0.10,0.15\}$), but the lift-error grows with $T$---from $8.6$ to
$31.8$ price\bp{} at $n=4$ as $T$ runs $0.05\to0.5$---so at long maturities the
lift still tops the solve floor at $n=16$ and the crossover slides to larger $n$.
Longer horizons simply need more factors.

\subsection{Bitcoin really is rough}
We put the model to real data. Calibrated to the live Deribit BTC surface with a
full-budget optimizer (differential evolution, $60$ generations plus local polish;
all $13$ maturities, up to $8$ points each, $97$ points), rough Heston returns
$H=0.090$, $\rho=-0.792$, $V_0=0.103$, $\theta=0.253$, $\lambda=2.567$,
$\nu=0.300$, at a vega-weighted implied-vol RMSE of $496$ vol\bp{}
(Figure~\ref{fig:calib}). The Hurst exponent lands firmly in the rough regime and
the correlation is strongly negative---both exactly what one expects of a crypto
market with a steep, persistent skew. We do not dress up the fit: at $496$ vol\bp{}
it is coarse, and the per-maturity smiles in Figure~\ref{fig:calib}(b) show where
it strains. That coarseness is a statement about the model's expressiveness on
wide, richly curved crypto smiles, not about our solver, and it is orthogonal to
the solve-error we have been tracking.

\begin{figure}[t]\centering
\includegraphics[width=0.98\linewidth]{fig2_calibration_fit.png}
\caption{Rough Heston calibrated to the Deribit BTC surface ($97$ points, $13$
maturities) at the canonical fit ($H=0.090$, $\rho=-0.79$; RMSE $496$ vol\bp).
(a)~market versus model implied volatility at all points (dashed line: identity;
colour $=$ maturity). (b)~per-maturity smiles for four representative maturities:
market (markers) versus model (dashed). The fit is deliberately coarse: the
residual spread reflects a known limitation of the lifted rough-Heston form on
wide, richly curved crypto smiles---a model-expressiveness issue, not a solver
artifact---and is orthogonal to the solve-error studied elsewhere.}
\label{fig:calib}
\end{figure}

\subsection{The parametric network, and the hybrid that makes it pay off}
\label{sec:hybrid}
A single network stretched over a five-parameter box cannot match a network
trained at one point, and we quantify the price of coverage. The parametric
network varies $(H,\nu,\rho,\lambda,\theta)$ while holding $V_0$ and $T$ fixed; at
the box centre its solve-error is $11.0\pm4.7$ price\bp{} against $5.9\pm1.2$ for a
single-point network at the same $n=10$---roughly a two-fold penalty. And the
penalty is uneven: it hovers near $9$--$11$ price\bp{} at typical points but
climbs to $47.7$ in the rough, low-$H$ corner and $29.8$ at high vol-of-vol, where
the PDE is stiffest. We report this plainly; a larger training budget would narrow
the corner gap.

The three calibration routes tell a nuanced story, and we report each result
plainly (Table~\ref{tab:calib}, Figure~\ref{fig:hybrid}). Fourier-in-the-loop
reaches $23.7$ vol\bp{}, at a cost of $2018$ Fourier evaluations and about
$168$ s. The parametric PINN on its own is almost free---about $2$ s, zero
Fourier calls---yet its fit is poor, $222\pm20$ vol\bp{} under the true model,
both because it calibrates to its own imperfect prices and because this
single-slice optimum happens to fall \emph{outside} the training box (it prefers
$\nu\approx0.87$, $\theta\approx0.39$). We therefore would not recommend the
standalone network as a calibrator. It does, however, make an effective starting
point: seeding the local Fourier polish with the PINN's guess reaches
$29.4\pm7.1$ vol\bp{}---on par with from-scratch on a $34$--$52\%$ smile---in only
$237\pm46$ Fourier evaluations and about $23$ s. That is roughly nine times fewer
true-model calls and seven times faster for the same practical fit, so the
parametric network is most useful not as a pricer but as a differentiable warm
start.

\begin{table}[t]\centering
\caption{Calibration of the BTC $T\approx0.15$ slice, three ways (3 PINN seeds;
all judged by the true Fourier model over the same feasible region).}
\label{tab:calib}
\begin{tabular}{lccc}
\toprule
method & fit RMSE (vol\bp) & wall-clock (s) & Fourier evals\\
\midrule
Fourier-loop from scratch      & $23.7$        & $168$ & $2018$\\
PINN-only (single-shot)        & $222\pm20$    & $\sim2$ & $0$\\
Hybrid (PINN $+$ Fourier polish) & $29.4\pm7.1$ & $23$  & $237\pm46$\\
\bottomrule
\end{tabular}
\end{table}

\begin{figure}[t]\centering
\includegraphics[width=0.82\linewidth]{fig3_hybrid.png}
\caption{Hybrid calibration. Fit quality (left) and Fourier-evaluation cost
(right) for the three methods on the BTC slice: the hybrid matches from-scratch
quality at $\sim\!9\times$ fewer Fourier evaluations, while the PINN alone is fast
but inaccurate.}
\label{fig:hybrid}
\end{figure}

% =====================================================================
\section{Discussion}

\paragraph{What the scaling result does and does not say.}
The finding we care about most is that the network's own error in solving the
lifted PDE degrades only gently as the dimension grows, while grid methods become
impossible. We have taken pains to state it in its honest, weak form. The
solve-error is not constant in $n$---a claim an early single-seed run seemed to
support and three-seed statistics refuted---it grows sub-linearly, roughly
doubling from $n=4$ to $n=32$. And $50^{n+1}$ is the naive grid; sparse-grid, ADI,
and tensor-train methods reach higher dimensions, so the defensible claim is that
a competitive finite-difference solve for this PDE at $n=16$--$32$ is impractical,
not impossible. Even so, the contrast---sub-linear error and bounded wall-clock
against a grid that cannot be stored---is a concrete instance of the advantage
that motivates mesh-free neural solvers, and it is the cleanest such instance we
know of on an economically real PDE.

\paragraph{A lesson beyond rough volatility.}
The most transferable idea here is methodological. When a surrogate replaces an
exact solver, it is dangerously easy to validate it against the wrong reference.
Our network solves the finite-$n$ lifted PDE, whose exact solution is the lifted
Monte Carlo, not the $n\to\infty$ model; comparing straight to the true model
conflates the network's solve-error with the model's lift-error. We found the two
have opposite sign and that the historical ``PINN-versus-truth'' gap was
lift-dominated, not the network's doing. Validate against a same-fidelity
reference, report the two errors separately, and let their crossover set the
approximation order: none of this is special to rough volatility, and we would
urge it on anyone placing a learned solver on top of an approximated operator.

\paragraph{The calibration story, straight.}
We are deliberately blunt about a negative result and its fix. One network spread
over five parameters is markedly less accurate than a point-trained one, worst in
the rough, high-vol-of-vol corners, and calibrating with it alone is fast but
unreliable---it fits its own imperfect prices, and the single-slice optimum can
sit outside the training box. Yet the same network is a fine \emph{warm start}:
seeding a short true-model polish with its guess recovers from-scratch quality at
about an order of magnitude fewer true-model evaluations. The value proposition is
amortized---train once, then recalibrate cheaply---with the exact pricer kept in
the loop for a brief, decisive finish.

\paragraph{Limitations.}
Our evidence rests on a single market snapshot (one Deribit BTC surface), and the
parametric and calibration studies use a single maturity slice; timings are
single-machine CPU figures at moderate training budgets. The parametric network
fixes $V_0$ and the maturity, its surface-level fit to the wide crypto smile is
coarse ($496$ vol\bp), and its solve-error carries a mild directional bias---a
slight over-pricing of at-the-money and call strikes that grows with $n$. We study
$n\le32$; the sub-linear trend need not continue arbitrarily far.

\paragraph{Where we take it next.}
Several extensions are immediate: add the maturity as a network input to price
whole surfaces rather than single slices; scale width, depth, and GPU training to
shrink the parametric corner penalty and probe larger $n$; carry the approach to
other rough models (rough Bergomi) and to American and path-dependent payoffs; and
map the warm-start polish across market regimes. We would also like to pin down
the high-$n$ solve-bias and whether loss reweighting or a residual correction
removes it, and to benchmark the mesh-free solver directly against structured
finite differences at the moderate $n$ where the latter still stand a chance.

\bibliographystyle{plainnat}
\bibliography{refs}
\end{document}
